\chapter{Experimental setup}
\section{OKTAVIAN}
\begin{figure}
  \centering
  \includegraphics[width=0.5\columnwidth]{eps/E-C.eps}
  \caption{simbol}
  \label{fig:simbol}
\end{figure}
\ref{fig:simbol}
\section{MAIKoのセットアップ}
\subsection{検出ガスの決定}
$\alpha$粒子のレンジで決めた。
efficiency が一番大きくなるように。
\subsection{ドリフトスピード}
ドリフト速度の決定方法は30 degree 方向に$\alpha$線源から$\alpha$を出して、
その飛跡がデータ上でどう見えるかで決定する。
ドリフト速度の時間依存性も見た。
\section{液体シンチレータのセットアップ}
